%% tab-platforms.tex -- Table: feature-by-feature against the prompt-lifecycle
%% platforms, which are the genuinely closest systems.
%%
%% This table exists because an earlier draft asserted that the release concern
%% was unoccupied. It is not. MLflow's Prompt Registry, LangSmith, and Langfuse
%% all ship versioned prompts, a named pointer the client resolves at call time,
%% rollback by re-pointing, and role restrictions on who may move the pointer.
%% A comparison that gave those systems "none" for governance was a straw man,
%% and the delta this paper can honestly claim is narrower than the one the
%% earlier draft claimed.
%%
%% Every row is sourced to dated product documentation. Product capability moves;
%% the access date is part of the claim, and the caption says so.

\begin{table}[tbp]
\centering
\caption{Feature-by-feature comparison with the prompt-lifecycle platforms --- the
closest systems to \sys, and the ones its contribution should be judged against.
\textbf{Claims are as of 2026-08-03}, sourced to the cited documentation; these are
live products and a reader checking later should expect drift. Above the rule are
the four concerns an earlier draft of this paper got wrong: versioning, call-time
pointer resolution, rollback, and role-restricted promotion are \emph{present} in
all three, and \sys is not novel in any of them. Below the rule is where a
difference survives scrutiny. \secref{sec:compare:platforms} argues both halves.}
\label{tab:platforms}
\small
\setlength{\tabcolsep}{4.5pt}
\renewcommand{\arraystretch}{1.10}
\begin{cbwide}
\begin{tabular}{@{}L{2.86cm}L{3.16cm}L{3.16cm}L{3.16cm}VL{3.30cm}@{}}
\toprule
\theadrow
\thead{Concern} &
\thead{MLflow Prompt Registry} & \thead{LangSmith} & \thead{Langfuse} &
\thead{\sysbare{}} \\
\midrule

\textbf{Prompt versioning} &
immutable commit-based versions, with diffs~\cite{mlflowpromptregistry} &
commit history per prompt~\cite{langsmithprompts} &
version identifiers per prompt~\cite{langfuseprompts} &
\cellcolor{cbControlBg}versions per \gasset{} --- the same mechanism, not a
better one \\

\zebra
\textbf{Call-time pointer} &
mutable \emph{aliases}~\cite{mlflowpromptregistry} &
commit tags; reserved \code{production}, pulled by name at
runtime~\cite{langsmithprompts} &
\emph{labels}; an unlabelled fetch serves \code{production}~\cite{langfuseprompts} &
\cellcolor{cbControlBg}\livetarget{} per family --- \textbf{not a differentiator}
\\

\textbf{Rollback} &
re-point the alias~\cite{mlflowpromptregistry} &
rollback to a prior commit~\cite{langsmithprompts} &
re-label an older version~\cite{langfuseprompts} &
\cellcolor{cbControlBg}restore a \emph{recorded} checkpoint --- governed path only
(\secref{sec:limitations}) \\

\zebra
\textbf{Who may promote} &
registry permissions, where the deployment provides them &
owners-only tag create and delete~\cite{langsmithprompts} &
\emph{protected labels}, by project role~\cite{langfuseprompts} &
\cellcolor{cbControlBg}\operator{} or \admin{} scope --- comparable, and no
separation of duties (\secref{sec:limitations}) \\

\midrule

\textbf{Families beyond prompts} &
models; no other artifact type is a registry citizen &
datasets &
datasets &
\cellcolor{cbControlBg}\statFamiliesW{} families --- workflows, tools, MCP
servers, corpora, skills --- under one mode vocabulary
(\tabref{tab:families}) \\

\zebra
\textbf{Evidence bound to the version} &
evaluation integrated; the score annotates &
experiments and datasets, linked &
evaluations and scores, linked &
\cellcolor{cbControlBg}scores, corpus identity, diagnosis, and optimizer choice
bound to the candidate and carried into review \\

\textbf{Score as a precondition on promotion} &
none documented: promotion is a pointer move &
none documented &
none documented &
\cellcolor{cbControlBg}a failing gate stops the candidate --- inside the
refinement state machine, and only there \\

\zebra
\textbf{Automated candidate proposal} &
\textsc{gepa} and metaprompting
optimizers~\cite{mlflowoptimize} &
none documented &
none documented &
\cellcolor{cbControlBg}optimizers wired \emph{to} the gate --- the optimization is
not the delta, the wiring is \\

\textbf{Cross-family impact} &
\chipnone{} --- no second family to impact &
\chipnone{} &
\chipnone{} &
\cellcolor{cbControlBg}a change is relatable to dependents, because dependents are
records \\

\bottomrule
\end{tabular}
\end{cbwide}
\tabnote{``None documented'' means the mechanism was not found in the cited pages
on the access date --- a claim about the documentation, which a reader can check,
rather than about the product, which we are not positioned to make. Tracing-only
systems such as Phoenix~\cite{phoenix} are excluded: they do not offer prompt
lifecycle management and these rows would score them unfairly.}
\end{table}
